\subsection{Ethical and organisational considerations}\label{sec:ethics}
A scheduler that monitors attention and fatigue and recommends when a person should work on what and when to rest touches on autonomy, privacy and fairness, and these issues must be designed in rather than added later \cite{parasuraman2000,feigh2012,amodei2016}. We see four requirements for responsible deployment. First, \emph{the system should be advisory}. Recommendations should be presented to the worker, who retains the authority to accept or override them; evidence from human--robot teaming indicates that the allocation of decision authority affects both efficiency and satisfaction \cite{gombolay2015}. Second, \emph{cognitive-state data are sensitive personal data}. Physiological and behavioural signals should be processed with informed consent, ideally on the worker's own device, and should never be used for performance appraisal. The observation model of AA-DQN needs only two scalar estimates per slot, which supports such data minimisation. Third, \emph{fairness across workers} requires that policies do not systematically disadvantage people whose fatigue dynamics differ from the nominal profile. Section~\ref{sec:robust} shows how performance changes for resilient and fatigue-prone profiles, and personalisation (Section~\ref{sec:future}) would further reduce such disparities. Fourth, \emph{safety constraints} should bound the policy, for example by guaranteeing a minimum amount of rest per period regardless of deadline pressure, consistent with safe RL principles \cite{garcia2015}. Notably, the learned policies in this study already increase rest relative to deadline-driven rules. The reward weight $\lambda$ makes the productivity--well-being trade-off explicit and auditable (Section~\ref{sec:pareto}), so stakeholders can choose it deliberately rather than leaving it implicit in a heuristic.

\subsection{Limitations}\label{sec:limitations}
The study has limitations that define its scope. (i) \emph{Simulation, not human data.} The cognitive model is grounded in established functional forms and its parameters were chosen to be behaviourally plausible, but it has not been calibrated against individual time-series of attention and fatigue. The absolute numbers should therefore be interpreted comparatively. The robustness analyses mitigate, but do not remove, this concern. (ii) \emph{Abstracted sensing.} Cognitive-state estimation is modelled as additive Gaussian noise on the true state. Real estimators can be biased, drift over the day, or fail intermittently, and they may be correlated with task type. (iii) \emph{A single worker and a single day.} Multi-day carry-over of fatigue, sleep, and team-level interactions such as hand-offs and shared deadlines are outside the scope. (iv) \emph{Action abstraction.} The agent chooses a difficulty class and delegates the choice of task within a class to EDF. This keeps the policy compact and size-invariant but excludes policies that, for example, deliberately abandon a hopeless task. (v) \emph{Algorithmic scope.} We study value-based agents with a finite history. Recurrent, model-based or policy-gradient agents \cite{hausknecht2015,igl2018,schulman2017} might extract more from the observation sequence. (vi) \emph{Simulated time-on-task effects only.} Motivation, task interest, interruptions and errors, which also shape real productivity \cite{mark2008,gonzalez2004,czerwinski2004}, are not modelled.

\subsection{Future work}\label{sec:future}
Three extensions follow directly. First, \emph{empirical calibration and validation}: fitting the latent dynamics to longitudinal data from vigilance probes, wearables and task logs \cite{basner2011,charles2019,shaffer2017}, then evaluating the scheduler in a controlled user study with NASA-TLX and performance outcomes \cite{hart1988}, ideally as a micro-randomised trial \cite{klasnja2015}. Second, \emph{personalisation}: domain randomisation over worker profiles during training \cite{tobin2017,zhao2020}, and Bayesian or meta-learning adaptation to an individual's parameters, in the spirit of personalised just-in-time interventions \cite{liao2020,nahumshani2018}. Third, \emph{multi-worker and human--robot settings}: extending the formulation to teams and to collaborative cells in which a robot absorbs demanding sub-tasks when a co-worker's attention is low \cite{ajoudani2018,lamon2019,zhang2022hrc,peternel2018}, with constrained or multi-objective RL to enforce explicit well-being guarantees \cite{hayes2022}.
